\chapter{Training Configuration and Domain Randomization}%
\label{app:training_config}%

\section{PPO Hyperparameters}%
\label{app:ppo_hyperparams}%

\Cref{tab:ppo_hyperparams} lists the \ac{PPO} hyperparameters used for teacher policy training.

\begin{table}[htb]
    \centering
    \caption{PPO hyperparameters for teacher training.}
    \label{tab:ppo_hyperparams}
    \begin{tabular}{lc}
        \toprule
        \textbf{Hyperparameter} & \textbf{Value} \\
        \midrule
        Learning rate & $1 \times 10^{-4}$ \\
        Clip parameter $\epsilon$ & 0.2 \\
        Mini-epochs per update & 4 \\
        Horizon length & 16 steps \\
        Minibatch size & 2048 \\
        Discount factor $\gamma$ & 0.99 \\
        GAE $\lambda$ & 0.95 \\
        Number of environments & 1{,}024 \\
        Max training epochs & 33{,}000 \\
        \bottomrule
    \end{tabular}
\end{table}

\section{Network Architecture}%
\label{app:network_architecture}%

The teacher policy uses an \ac{LSTM}-based recurrent network implemented in rl\_games.
The actor consists of an MLP encoder followed by an LSTM layer with 1024 hidden units, followed by an output head mapping to the 23-dimensional action space.
The critic has an identical architecture but receives the larger critic observation (\Cref{tab:ac_observations}).

The student uses a vision--proprioception fusion architecture:
\begin{enumerate}
    \item \textbf{Visual encoder:} A ResNet-18 backbone processes each stereo image independently, producing a 32-dimensional feature vector per image.
    \item \textbf{Cross-attention fusion:} A cross-attention transformer module fuses the two stereo features into a single 64-dimensional visual embedding.
    \item \textbf{Policy head:} The visual embedding is concatenated with the proprioceptive state vector and fed into an \ac{LSTM} with 512 hidden units, followed by a linear output layer producing the 23-dimensional action.
\end{enumerate}

\section{Automatic Domain Randomization Parameters}%
\label{app:adr_params}%

\Cref{tab:adr_params} lists the key randomized parameters and their start/max ranges.
The maximum ranges for finger joint stiffness and damping correspond to values measured on the real Agile Hand hardware using an Isaac Lab gain tuner.
The thumb rotation velocity limit is ramped down toward the real hardware limit of approximately 12\,deg/s ($\approx 0.21$\,rad/s).
All \ac{ADR} parameters share a single increment counter ($N_{\text{max}} = 50$, success threshold $\tau_{\text{ADR}} = 0.4$), meaning all parameters progress together.

\begin{table}[htb]
    \centering
    \caption{Automatic Domain Randomization parameters.}
    \label{tab:adr_params}
    \begin{tabular}{lcc}
        \toprule
        \textbf{Parameter} & \textbf{Start Range} & \textbf{Max Range} \\
        \midrule
        Object mass scale & $[1.0, 1.0]$ & $[0.5, 2.0]$ \\
        Object friction (static) & $[1.0, 1.0]$ & $[0.5, 1.5]$ \\
        Robot restitution & $[1.0, 1.0]$ & $[0.8, 1.0]$ \\
        Arm stiffness scale & $[1.0, 1.0]$ & $[0.5, 2.0]$ \\
        Arm damping scale & $[1.0, 1.0]$ & $[0.5, 2.0]$ \\
        Finger MCP pitch stiffness & $[1.0, 1.0]$ & $[0.17, 1.0]$ \\
        Finger MCP pitch damping & $[1.0, 1.0]$ & $[0.083, 1.0]$ \\
        Finger MCP yaw stiffness & $[1.0, 1.0]$ & $[0.025, 1.0]$ \\
        Finger PIP stiffness & $[1.0, 1.0]$ & $[0.02, 1.0]$ \\
        Thumb rotation velocity limit & $[20.0, 20.0]$ & $[0.21, 20.0]$\,rad/s \\
        \bottomrule
    \end{tabular}
\end{table}

\begin{figure}[htb]
    \centering
    \includegraphics[width=0.75\textwidth]{plots/adr_ranges}
    \caption{ADR parameter ranges. Blue bars indicate the full randomization range at maximum ADR level. Red diamonds mark the nominal start values. The thumb velocity limit spans the widest relative range, from 20\,rad/s (simulation) to 0.21\,rad/s (hardware limit).}
    \label{fig:adr_ranges}
\end{figure}

\section{Full Results Including Physics Instability}%
\label{app:full_results}%

The results in \Cref{ch:experiments} exclude physics instability events (NaN joint states, velocity explosions), which are simulation-only artifacts that cannot occur on physical hardware.
This section reports the full evaluation data including physics instability for completeness.

\begin{table}[htb]
    \centering
    \caption{Teacher policy per-object evaluation including physics instability (200 episodes per object, ADR-0 physics).}
    \label{tab:teacher_eval_full}
    \small
    \begin{tabular}{lrrrrrr}
        \toprule
        \textbf{Object} & \textbf{Lift (\%)} & \textbf{UER (\%)} & \textbf{Coll. (\%)} & \textbf{OOB (\%)} & \textbf{Palm (\%)} & \textbf{Phys. (\%)} \\
        \midrule
        basketball shoe       &  99.0 & 32.0 &  7.0 & 21.0 &  0.0 &  4.0 \\
        closed fist           & 100.0 & 19.5 &  0.0 &  9.0 &  0.0 & 10.5 \\
        elephant toy          & 100.0 & 22.0 &  2.0 &  0.5 &  0.5 & 19.0 \\
        mario                 &  98.0 &  8.5 &  0.0 &  1.0 &  0.0 &  6.5 \\
        milk pot              &  96.5 & 15.0 &  0.5 &  5.0 &  0.0 &  9.5 \\
        teddy bear            &  99.5 & 21.5 &  0.5 & 14.5 &  0.0 &  6.5 \\
        toy excavator         &  99.0 & 25.5 &  2.0 &  7.5 &  1.0 & 15.0 \\
        turtle candle holder  & 100.0 & 35.0 &  0.0 & 24.0 &  0.0 & 10.5 \\
        \midrule
        \textbf{Average (8 objects)} & \textbf{99.0} & \textbf{22.4} & \textbf{1.5} & \textbf{10.3} & \textbf{0.2} & \textbf{10.2} \\
        \bottomrule
    \end{tabular}
\end{table}

\section{Per-Object Failure Mode Breakdowns}%
\label{app:per_object_failure}%

\Cref{fig:per_object_unsafe_heatmap,fig:per_object_unsafe_heatmap_physics} show the per-object unsafe-episode rate across all swept SafeDAgger thresholds, with \ac{BC} ($\delta \to 0$) and DAgger ($\delta \to \infty$) as endpoints.
\Cref{fig:per_object_unsafe_heatmap} excludes physics-instability events; \Cref{fig:per_object_unsafe_heatmap_physics} includes them for completeness.

\begin{figure}[htbp]
    \centering
    \includegraphics[width=\textwidth]{plots/run14/run14_per_object_unsafe_heatmap}
    \caption{Per-object unsafe-episode rate across the threshold sweep (physics excluded). Each cell shows the \acs{UER} for one object at one threshold. The basketball shoe and turtle candle holder exhibit the highest rates across most thresholds.}
    \label{fig:per_object_unsafe_heatmap}
\end{figure}

\begin{figure}[htbp]
    \centering
    \includegraphics[width=\textwidth]{plots/run14/run14_per_object_unsafe_heatmap_with_physics}
    \caption{Per-object unsafe-episode rate across the threshold sweep (physics included). Physics-instability events raise the overall rates but do not change the relative ordering across methods.}
    \label{fig:per_object_unsafe_heatmap_physics}
\end{figure}

%
